\documentclass[12pt]{article}
\usepackage[utf8]{inputenc}
\usepackage[T1]{fontenc}
\usepackage{hyperref}
\title{Southern Muji Language Resources for K-12 Curriculum Development}
\author{Compiled by Claude Sonnet 5.5 (claude-sonnet-5-5) for Andr\'as Kornai}
\date{}
\begin{document}
\maketitle

\section{Introduction}
Southern Muji (ISO 639-3: ymc; Glottocode sout2722) is a Sino-Tibetan (Ngwi/Loloish) language of southern Yunnan, China. The largest concentrations of speakers are in Gejiu and Mengzi counties, followed by Hekou, Pingbian and Jinping. Estimates of speaker numbers differ by source: Verbix gives 26,000 and Joshua Project about 29,000. The speakers are officially counted within the Yi nationality. The code ymc was split out of the older umbrella code ymj (Muji Yi) in 2008, alongside Mojii, Qila Muji, Northern Muji and Muzi, so sources sometimes conflate these languages. Global Recordings Network lists the varieties Dazhai, Gamadii, Jinhe, Ma'andi, Tongchang and Yingpan.

The searches found no primers, textbooks or school programmes for Southern Muji, and Joshua Project notes that Bible translation is still needed.

\section{Primary Educational Resources}
No Southern Muji primers or textbooks turned up in the searches. The language is not known to be taught in schools.

\section{Linguistic and Reference Materials}
\begin{itemize}
\item Wikipedia, Moji language (related Muji varieties): \url{https://en.wikipedia.org/wiki/Moji_language}
\item Wikipedia, Muzi language (a separate Muji variety, not ymc): \url{https://en.wikipedia.org/wiki/Muzi_language}
\item Wikipedia, Southeastern Loloish languages: \url{https://en.wikipedia.org/wiki/Southeastern_Loloish_languages}
\item Verbix Endangered Languages, Southern Muji: \url{https://docs.verbix.com/EndangeredLanguages/SouthernMuji}
\end{itemize}

\section{Religious and Community Resources}
\begin{itemize}
\item Joshua Project, Southern Muji in China: \url{https://joshuaproject.net/people_groups/18613/CH}
\item Asia Harvest, Muji people group profile (PDF): \url{https://people-groups.asiaharvest.org/people-group-profiles/yunnan/muji.pdf}
\end{itemize}

\section{Digital and Audio Resources}
\begin{itemize}
\item Global Recordings Network page for Muji, Southern (ymc). The page states that no recordings are currently available in this language: \url{https://globalrecordings.net/en/language/ymc}
\end{itemize}

\section{Curriculum Development Context}
Southern Muji is reported as vulnerable but widely used as a first language in its community. It has no known written standard or school materials, so a curriculum would start with orthography design and community-led documentation. Any work should take care not to mix Southern Muji with the other Muji languages that share the name.

\end{document}